% Complete root derivation, prefix HZD. No existence of an off-real zero is assumed.
\section{The supported heat divisor and its exact arithmetic and lattice fibres}
\label{sec:full-lattice-heat-divisor}

\paragraph{Scope and sources.}
The original heat family and all its constants are those of Brad Rodgers
and Terence Tao, \href{https://arxiv.org/abs/1801.05914v5}{\emph{The
De Bruijn--Newman constant is non-negative}}, original
\nolinkurl{fmp-template.tex}, equations \texttt{hoz}, \texttt{sas},
\texttt{phidef} and \texttt{htdef}, lines 78--103.
The full coefficient carrier comes from The Clankers,
\href{https://github.com/KokunoYumeto/zeta-function-research-reader/blob/a2851f9eb352e7e4376bc629de86fa4ed9c1c434/workbenches/splitzero-tandem/foundations/split-support-geometry-v11/split_support_geometry_arithmetic_curve_v11.tex#L1226}{\emph{Split
Support Geometry}, v11, \texttt{def:lattice-split}}.
The exact characteristic-one extension is ASL18--22 in this reader,
with its cited Connes--Consani sources. All quotient and local-block
comparisons used in this section are proved below. In particular a
supported-zero equation is not silently changed into an absorbing-zero
equation. No off-real zero is postulated.

\subsection{The supported equation and the full coefficient square}

Let $R$ be a commutative unital ring, allowing the zero ring, and let
$L$ be the entire nontrivial bounded distributive lattice. Write
\begin{equation}\tag{HZD1}
 S_R=G_L(R),\qquad \widehat f=(f,1_L),\qquad
 e_R=(0_R,1_L),\qquad \tau_R=(0_R,0_L).
\end{equation}
Its operations are those of ZIL1. In the zero-ring case this is
exactly $L$, with $e_R=1_{S_R}$ and $\tau_R=0_{S_R}$.
For a given $f\in R$, the supported equation $\widehat f=e_R$
defines the congruence quotient
\begin{equation}\tag{HZD2}
 \mathcal D_L(R,f):=S_R/\langle\widehat f\sim e_R\rangle
       \simeq G_L(R/(f)),\qquad
 (r,\lambda)\longmapsto([r],\lambda).
\end{equation}
Here angle brackets denote the smallest semiring congruence containing
the indicated pair; no ideal congruence is substituted for it.

To prove (HZD2), the displayed map is well defined on the full carrier,
preserves both operations and both identities, and is surjective.
It identifies $\widehat f$ with $e_R$. Two elements in its kernel
congruence have the same label and amplitudes differing by a multiple
of $f$. Below top support both amplitudes must be zero, so those
elements were already identical. At top support write $r-s=af$.
Multiplying the generating relation by $\widehat a$ gives
$\widehat{af}\sim e_R$; adding $\widehat s$ gives
$\widehat r\sim\widehat s$, since $e_R+\widehat s=\widehat s$.
This proves that its whole kernel congruence is generated by the
stated pair, proving the isomorphism without deleting any label.

Let $K$ be the original tropical semifield with zero,
$K_L=K\otimes_{\mathbb B}L$, and
$\iota_L(\lambda)=1_K\otimes\lambda$. Keep the proved retraction
$\rho_L:K_L\to L$, $\rho_L\iota_L=\operatorname{id}_L$ of ZIL11.
The complete map is
\begin{equation}\tag{HZD3}
 \Phi_R:S_R\hookrightarrow R\times K_L,
       \qquad(r,\lambda)\longmapsto(r,\iota_L(\lambda)).
\end{equation}
It is injective, including for the zero ring: equality of ring
coordinates and application of $\rho_L$ give equality of both
source coordinates. Its image keeps the original constraint
$r\ne0_R\Rightarrow\lambda=1_L$.

The supported divisor commutes with this complete map in the exact
square
\begin{equation}\tag{HZD4}
\begin{CD}
 S_R @>{\Phi_R}>> R\times K_L\\
 @VV{(r,\lambda)\mapsto([r],\lambda)}V
 @VV{(r,w)\mapsto([r],w)}V\\
 G_L(R/(f)) @>{\Phi_{R/(f)}}>> (R/(f))\times K_L.
\end{CD}
\end{equation}
In the upper-right semiring the congruence generated by
$(f,1_{K_L})\sim(0_R,1_{K_L})$ is exactly equality of the second
coordinate together with ring-coordinate equality modulo $(f)$.
Indeed multiply the generating relation by $(a,0_{K_L})$ and add
$(s,w)$ to obtain $(s+af,w)\sim(s,w)$. Conversely reduction modulo
$(f)$ preserves the generating relation and distinguishes all remaining
coordinates. This proves the claimed congruence exactly. It also
proves that (HZD4) is a pushout: compatible maps out of its upper
and left objects impose precisely that one generating relation on
the upper-right object, and therefore factor uniquely through the
displayed lower-right quotient. The lower injection shows that the
supported divisor itself is retained faithfully there.

The different equation $\widehat f=\tau_R$ has the different quotient
\begin{equation}\tag{HZD5}
 S_R/\langle\widehat f\sim\tau_R\rangle\simeq R/(f).
\end{equation}
To prove this, multiply its generating relation by $e_R$. Since
$\widehat f e_R=e_R$ and $\tau_Re_R=\tau_R$, it imposes
$e_R\sim\tau_R$. This identifies every $z_\lambda$ with $\tau_R$
by multiplication, and gives the ring reflection $S_R\to R$.
The remaining relation is $f=0_R$, giving $R/(f)$. Conversely the
map $(r,\lambda)\mapsto[r]$ preserves that relation and is surjective,
so there are no additional identifications. Thus the exact connecting
map from (HZD2) to (HZD5) is the amplitude map
$G_L(R/(f))\to R/(f)$. This proves the relation between the two
zero conventions instead of replacing one with the other.

Inverting supported zero in the actual supported divisor gives
\begin{equation}\tag{HZD6}
 \mathcal D_L(R,f)[e^{-1}]\simeq L,\qquad
 ([r],\lambda)\longmapsto\lambda.
\end{equation}
This is ZIL3 applied to $e$, whose corner consists of all labelled
zeros and has identity $e$. The same conclusion follows by first
localizing $S_R$ to $L$: the supported equation already holds there
because both $\widehat f$ and $e_R$ map to $1_L$. Thus localization
commutes with the supported divisor, but its output is independent
of $f$. It does not erase the distinction between (HZD2) and (HZD5)
before localization.

\subsection{The single off-real heat-divisor object}

Use the original complex coordinate $z$, and set
$\Omega=\mathbb C\setminus\mathbb R$. Let
$A=\mathcal O(\Omega)$ be the unital ring of holomorphic functions
on this entire open set, with both connected components retained.
For the actual heat family define, for every real $t$,
\begin{equation}\tag{HZD7}
 R_t=A/(H_t|_\Omega),\qquad
 \mathscr D_{t,L}=G_L(R_t),\qquad
 \Phi_{R_t}:\mathscr D_{t,L}\hookrightarrow R_t\times K_L.
\end{equation}
This is the explicit quotient (HZD2) for the original heat function,
with its quotient map, full carrier, joins, meets and constrained
faithful image. The counterfactual object at the original time is
$\mathscr D_{0,L}$, together with the distinguished element $e$ and
the map (HZD7). It has been constructed without assuming that it
has an off-real arithmetic point.

For any holomorphic function $h$ on $\Omega$, the following
equivalences hold:
\begin{equation}\tag{HZD8}
 \mathcal O(\Omega)/(h)=0
 \ \Longleftrightarrow\ h\text{ is a unit in }\mathcal O(\Omega)
 \ \Longleftrightarrow\ h\text{ has no zero in }\Omega.
\end{equation}
The first equivalence says exactly that $1$ lies in the principal
ideal $(h)$. A unit cannot vanish, by evaluating an inverse identity
at a proposed zero. If $h$ has no zero, the function $1/h$ is defined
at every point of $\Omega$ and is holomorphic there. For completeness,
near any point $z_0$ write $h(z)=h(z_0)(1+v(z))$ with $v(z_0)=0$;
on a sufficiently small neighbourhood $|v|<1$, and the convergent
geometric series $h(z_0)^{-1}\sum_{n\ge0}(-v(z))^n$ is the local
holomorphic reciprocal. These local reciprocals agree on overlaps,
giving the global inverse on both components. This proves (HZD8).

The canonical support map on (HZD7) consequently satisfies
\begin{equation}\tag{HZD9}
\begin{gathered}
 \sigma_t:\mathscr D_{t,L}\longrightarrow L\text{ is an isomorphism}
 \quad\Longleftrightarrow\quad R_t=0\\
 \quad\Longleftrightarrow\quad
 H_t\text{ has no off-real zero}.
\end{gathered}
\end{equation}
If $R_t=0$, the carrier is precisely $L$ and $\sigma_t$ is its
identity. If $R_t\ne0$, the distinct elements $(0,1_L)$ and
$(1,1_L)$ have the same support; hence the specified map is not
injective. Together with (HZD8), this proves every direction of
(HZD9). The statement is about this canonical map, not an
unsupported abstract identification of two semirings.

The actual arithmetic part of its spectrum is
\begin{equation}\tag{HZD10}
 V_{\mathscr D_{t,L}}(e)\simeq\operatorname{Spec}R_t,
 \qquad
 D_{\mathscr D_{t,L}}(e)\simeq\operatorname{SpecLat}L.
\end{equation}
The explicit primes and their inclusions are exactly ZIL9, whose
proof applies to every nonzero $R_t$. If $R_t=0$, then $e=1$,
the first displayed space is empty and the second is the whole
lattice spectrum. Each actual off-real root $z_0$ gives the
surjective evaluation $R_t\to\mathbb C$ and the prime
\begin{equation}\tag{HZD11}
 Q_{z_0}=\{z_\lambda:\lambda\in L\}
       \cup\{([h],1_L):h(z_0)=0\}.
\end{equation}
Evaluation is well defined because $H_t(z_0)=0$, and its kernel
is prime because its target is a field. Conversely the absence
of every such root forces $R_t=0$ by (HZD8), so there is then
no arithmetic prime of any kind. We have not assumed that every
maximal ideal of a ring of holomorphic functions is an evaluation
ideal. No such assertion is needed for this exact emptiness test.

At $t=0$ the unchanged Rodgers--Tao identity gives
\begin{equation}\tag{HZD12}
 H_0(z)=\frac18\xi\left(\frac12+\frac{iz}{2}\right),\qquad
 s=\frac12+\frac{iz}{2},\quad z=-2i(s-\tfrac12).
\end{equation}
For $z=x+iy$, $\operatorname{Re}s=(1-y)/2$; therefore off-real
$z$ corresponds exactly to off-critical $s$. This is the source's
unchanged nontrivial-zero coordinate map. Thus the full supported
counterfactual object has an arithmetic root exactly when the
original zeta function has a counterexample. Supported-zero
invertibility in its localization is already proved by (HZD6),
and occurs in both cases of (HZD9). Whether $\sigma_0$ is an
isomorphism has not been decided by that localization calculation.

\subsection{Every actual primary root retains its complete supported jet}

The following construction applies to each actual root, without
postulating that an off-real one exists. Fix a real $t$ and a root
$z_0$ of the actual nonzero entire function $H_t$. Let $m$ be its
exact positive multiplicity. Retain the actual local factor
\begin{equation}\tag{HZD13}
 H_t(z)=(z-z_0)^m u(z),\qquad
 u(z)=\sum_{j\ge0}\frac{H_t^{(m+j)}(z_0)}{(m+j)!}(z-z_0)^j,
 \quad u(z_0)=\frac{H_t^{(m)}(z_0)}{m!}\ne0.
\end{equation}
This follows by the original Taylor expansion at $z_0$ and the
definition of the exact multiplicity. No coefficient or unit
factor has been reset to one.

Write $B=\mathcal O(\mathbb C)$ and
$\mathfrak m_{z_0}=\{h:h(z_0)=0\}$. The exact jet map is
\begin{equation}\tag{HZD14}
\begin{gathered}
 B_{\mathfrak m_{z_0}}/(H_t)
   \xrightarrow{\ \sim\ }\mathbb C[\epsilon]/(\epsilon^m),\\
 h\longmapsto\sum_{j=0}^{m-1}\frac{h^{(j)}(z_0)}{j!}\epsilon^j,
 \qquad z\longmapsto z_0+\epsilon .
\end{gathered}
\end{equation}
For fractions the map is extended by the inverse of the denominator
jet, whose constant term is nonzero. It is a ring homomorphism by
the product rule with its exact binomial coefficients. It is
surjective because every truncated polynomial is represented by
that polynomial in $z-z_0$. Its kernel consists of functions
vanishing to order at least $m$. Such an entire numerator $h$
has $h/(z-z_0)^m$ entire, by its Taylor expansion at $z_0$ and
ordinary division away from $z_0$. The entire function
$u=H_t/(z-z_0)^m$ has nonzero value at $z_0$, so is inverted in
$B_{\mathfrak m_{z_0}}$. Thus this numerator is a multiple of
$H_t$ there. The converse follows from (HZD13), proving the
kernel assertion and the isomorphism, while keeping the original
factor $u$ explicitly available.

For an off-real $z_0$, restriction $B\to A=\mathcal O(\Omega)$
induces the exact comparison
\[
 B_{\mathfrak m_{z_0}}/(H_t)
   \longrightarrow A_{\mathfrak m_{z_0}}/(H_t|_\Omega).
\]
The same Taylor proof applies on $\Omega$: division by
$(z-z_0)^m$ extends holomorphically at $z_0$ and is already
holomorphic at every other point of both components. Both
quotients therefore map isomorphically to the displayed jet
ring, and restriction preserves all derivatives at $z_0$.
The comparison is consequently an isomorphism, identifying
this block with the actual arithmetic local block of (HZD7).

Applying the exact supported quotient and localization maps gives
\begin{equation}\tag{HZD15}
 G_L(B_{\mathfrak m_{z_0}}/(H_t))
   \simeq G_L(\mathbb C[\epsilon]/(\epsilon^m)).
\end{equation}
The map is the coefficient map (HZD14) with every label unchanged;
its inverse is the inverse ring map with the same label. To check
localization directly, a top-supported denominator $(a,1_L)$
with $a(z_0)\ne0$ sends $(h,\lambda)$ to $(h/a,\lambda)$.
The semiring fraction relation is precisely ring-fraction equality
and equality of the labels: multiplying by any top-supported
denominator preserves that label. Thus no label is lost in the
localization used in (HZD15).

The jet object in (HZD15) has the exact coefficient realization
\begin{equation}\tag{HZD16}
 \left(\sum_{j=0}^{m-1}a_j\epsilon^j,\lambda\right)
       \longmapsto((a_0,\lambda),\ldots,(a_{m-1},\lambda)).
\end{equation}
Its image has all coordinate labels equal, and a nonzero
coefficient forces that common label to be top. This is a
supported submodule of $G_L(\mathbb C)^m$, not an identification
with that entire independently labelled free module. Let $E_m$
be the matrix with amplitude identity and every entry at top
support. Its action on the free module keeps each amplitude
and replaces every coordinate label by the join of all labels.
Thus $E_m^2=E_m$, and its image is exactly (HZD16). This also
constructs the split projective realization and its entire
projection without changing the coefficient basis.
The algebra product on this image is the retained truncated
Cauchy product, not coordinatewise multiplication:
\[
 ((a_j,\lambda))_j\,((b_j,\nu))_j
   =\left(\left(\sum_{r+s=j}a_rb_s,\lambda\wedge\nu\right)\right)_{0\le j<m}.
\]
Every coefficient sum has at least one term and every term has
the same label $\lambda\wedge\nu$, so this formula follows
directly from the original semiring operations even when its
amplitude cancels. Thus (HZD16) also preserves the entire jet
algebra on its stated image.

Multiplication by the original local coordinate $\epsilon=z-z_0$
is represented on that image by the matrix $N_m$ whose amplitude
has entries $1$ at $i=j+1$ and $0$ elsewhere, with \emph{every}
matrix entry at top support. Direct coefficient multiplication
gives
\begin{equation}\tag{HZD17}
\begin{gathered}
 N_m^m=Z_m,\qquad
 Z_m=((0,1_L))_{0\le i,j<m},\\
 Z_m((a_j,\lambda)_j)=((0,\lambda))_j,\qquad
 Z_m^2=Z_m,\qquad N_m^0=E_m
 \quad\text{on the projective image}.
\end{gathered}
\end{equation}
Indeed each application shifts the amplitude vector by one place,
with zero in the first position; after $m$ applications all
amplitudes vanish. All terms in every matrix product carry top
support, so the resulting entries are $e$, not $\tau$.
The stated action and idempotence follow by joining the equal
input labels. The categorical zero map instead has output
$(\tau,\ldots,\tau)$ and is different whenever the input label
is nonzero. This is the exact supported nilpotence, with its
actual null endomorphism retained.

Finally, supported-zero inversion sends the jet algebra to $L$.
The top-supported element $\widehat\epsilon$ and its power
$\widehat\epsilon^{\,m}=e$ all map to $1_L$:
\begin{equation}\tag{HZD18}
 G_L(\mathbb C[\epsilon]/(\epsilon^m))[e^{-1}]
       \simeq L,\qquad
 \widehat\epsilon\longmapsto1_L,\quad e\longmapsto1_L.
\end{equation}
Thus the amplitude-nilpotent coordinate becomes the identity on
the localized support module. In the projective realization,
both $N_m$ and $E_m$ become the all-top join matrix; its image
consists of constant label vectors and is canonically $L$.
This proves the same conclusion using the complete matrix map,
not just the scalar quotient. Before localization the map
$\Phi_{\mathbb C[\epsilon]/(\epsilon^m)}$ retains the entire
jet amplitude, its multiplicity, the original coordinate
$z_0+\epsilon$, and every support label. The localization
does not determine the existence or position of the root whose
local object has just been calculated.
